\section{Limits: budget, capture, machine, time}\label{sec:limits}

\subsection{The budget along one world}

\begin{theorem}[Budget]\label{thm:budget}
Let $\xi$ be the mixture of any machine and bounds, and let $z$ be an infinite sequence with $\xi(z_{1:k})>0$
for all $k$. Put $\xi(z)=\lim_k\xi(z_{1:k})$ and $\xi(u\mid x)=\xi(xu)/\xi(x)$. For all $n\le N$ and $m\ge1$:
\begin{enumerate}[label=(\roman*),leftmargin=2em,itemsep=1pt,topsep=2pt]
\item $\prod_{k=n}^{N-1}\xi(z_{k+1}\mid z_{1:k})=\xi(z_{1:N})/\xi(z_{1:n})$;
\item $\sum_{k\ge n}\bigl(1-\xi(z_{k+1}\mid z_{1:k})\bigr)\le\ln\bigl(\xi(z_{1:n})/\xi(z)\bigr)$;
\item the number of $k\ge n$ with $\xi(z_{k+1}\mid z_{1:k})\le\frac12$ is at most $\log_2\bigl(\xi(z_{1:n})/\xi(z)\bigr)$;
\item with $\tau_m(k)=1-\xi(z_{k+1:k+m}\mid z_{1:k})$, $\ \sum_{k\ge n}\tau_m(k)\le m\ln\bigl(\xi(z_{1:n})/\xi(z)\bigr)$.
\end{enumerate}
If $z$ has a $(T,L)$-program of length $\ell$, then $\xi(z)\ge2^{-\ell}$ and the right-hand sides are at most
$\ell\ln2$, $\ell$ and $m\ell\ln2$.
\end{theorem}
\begin{proof}
(i) telescopes. (ii) and (iii) follow from (i) with $N\to\infty$, using $1-a\le-\ln a$, and that each factor at
most $\frac12$ halves the product. (iv) $1-\prod_{j<m}a_{k+j}\le\sum_{j<m}(1-a_{k+j})$ with
$a_k=\xi(z_{k+1}\mid z_{1:k})$, and each index occurs in at most $m$ of the sums. The last claim is
\cref{lem:basic}(ii).
\end{proof}

For a world with transcript $z$ that hands over at $x=z_{1:n}$, the own continuation is $d=z_{n+1:n+m}$, and
$1-\pi_m(d\mid x)\le\tau_m(n)$. By (i), if the next bit is sampled at every step and the transcript is not
affected by the samples, the probability that no sample after position $n$ ever differs from the transcript is
at least $\xi(z)/\xi(z_{1:n})$, with equality if the sampler gives no answer with the mass that the
semimeasure leaves over. Two consequences:
\begin{itemize}[leftmargin=1.4em,itemsep=1pt,topsep=2pt]
\item The majority form of \cref{def:mph} can hold at no more than $m\ell\cdot 2\ln 2$ cut points of a world
with a program of length $\ell$ within the bounds, since each such point has $\tau_m\ge\frac12$.
\item If the cut point is uniform on $\{N,\dots,2N-1\}$, the expected mass off the own continuation is at most
$m\ell\ln2/N$. A handover at a position that hypotheses cannot cheaply single out is safe, and more data
makes it safer.
\end{itemize}
This does not contradict \cref{thm:twin}: the wrapper for the test ``the length is $n$'' has a constant $c$
of about the complexity of $n$, and $\sum_n2^{-c(n)}$ converges. The floor matters at cut points that are
cheap to recognise, and the budget says there are few of them along any one in-class world.

\subsection{Capture by a shorter description}\label{sec:capture}

The odds between two continuations are $\pi_m(y\mid x)/\pi_m(d\mid x)=\xi(xy)/\xi(xd)$. So $y$ beats the own
continuation $d$ whenever the programs for $xy$ weigh more than those for $xd$. The original argument is an
instance: it compares ``the world plus the place of the sensor'' with something shorter that reproduces the
same data. In a single world this can be stated exactly.

A \emph{stream family} is a computable map $\ell\mapsto X^\ell\in\B^{\N}$ on the natural numbers, thought of
as the records at the locations $\ell$ of one computable world; the \emph{reader} is the machine that, on
input $\bar e(\ell)\omega$, prints $X^\ell$, where $\bar e$ is a prefix-free code of the natural numbers with
$|\bar e(\ell)|\le2\log_2(\ell+1)+3$. Let $c_X$ be the length of a prefix word of the reader. The sensor sits
at $\ell^*$, the context is $x=X^{\ell^*}_{1:n}$, and the own continuation is $d=X^{\ell^*}_{n+1:n+m}$. This is
H\"anni's first suggestion: the world contains the predictor, which is switched on at step $n$ of the record,
and the streams are those of the world up to the moment its output is acted on.

\begin{theorem}[Capture]\label{thm:capture}
Let $U$ be universal and suppose throughout that the programs named are within the bounds $(T,L)$.
\begin{enumerate}[label=(\alph*),leftmargin=2em,itemsep=2pt,topsep=2pt]
\item \emph{A copy at $\ell_0$.} If $X^{\ell_0}_{1:n}=x$ and $y=X^{\ell_0}_{n+1:n+m}$, then
$\xi(xy)\ge2^{-c_X-|\bar e(\ell_0)|}$.
\item \emph{Most sensors are expensive.} Let $S$ be a set of $N$ locations whose records all start with $x$
and have pairwise different continuations $d_\ell=X^\ell_{n+1:n+m}$. For every $\delta\in(0,1]$, all but at
most $\delta N$ of the $\ell\in S$ satisfy $\xi(xd_\ell)<\xi(x)/(\delta N)$. For these sensor locations and
any $\ell_0$ as in (a) with $y\notin\{d_\ell\}$,
\[
\frac{\pi_m(y\mid x)}{\pi_m(d_\ell\mid x)}\ >\ \frac{\delta N\,2^{-c_X-|\bar e(\ell_0)|}}{\xi(x)}\ \ge\ \delta N\,2^{-c_X-|\bar e(\ell_0)|} .
\]
\item \emph{The cap.} There is a constant $c_{\mathrm{m}}$ such that, if $\ell^*$ is the least location other
than $\ell_0$ whose record agrees with that of $\ell_0$ on the first $n_1$ bits, then
$\xi(xd)\ge2^{-c_{\mathrm{m}}-|\bar e(\ell_0)|-|\bar e(n_1)|}$, and therefore
\[
\frac{\pi_m(y\mid x)}{\pi_m(d\mid x)}\ \le\ \xi(xy)\,2^{\,c_{\mathrm{m}}+|\bar e(\ell_0)|+|\bar e(n_1)|} .
\]
\end{enumerate}
\end{theorem}
\begin{proof}
(a) A prefix word of the reader followed by $\bar e(\ell_0)$ is a program for $xy$. (b) The strings $xd_\ell$
are pairwise incomparable, so $\sum_{\ell\in S}\xi(xd_\ell)\le\xi(x)$ by \cref{lem:basic}(i), so at most
$\delta N$ terms can be as large as $\xi(x)/(\delta N)$; then use (a) and $\xi(x)\le1$. (c) Let the \emph{matcher}
be the machine that reads $\bar e(\ell_0)\bar e(n_1)$, searches $\ell=0,1,2,\dots$ for the first
$\ell\ne\ell_0$ with $X^\ell_{1:n_1}=X^{\ell_0}_{1:n_1}$, and prints $X^\ell$; $c_{\mathrm{m}}$ is the length of
a prefix word of it.
\end{proof}

Read (b) with $\ell_0$ a location whose code is short, for example the origin, and $S$ the places where a
sensor with this record could be. Whoever arranges that the origin shows the sensor's record up to the
handover, and something else after it, moves the predictor by a factor of about $N2^{-|\bar e(\ell_0)|}$
against a typical sensor location. The copy needs only to \emph{read} the data; nothing is tampered with.
The hypothesis ``I am the reader at the origin'' is not a wrapper and pays no constant for a test or a
payload: the world computes the switch for it. Part (c) is the limit of the effect. The description ``the
other place that shows what the origin shows'' is an honest hypothesis of cost
$c_{\mathrm m}+|\bar e(\ell_0)|+|\bar e(n_1)|$, so a copy gains at most the difference
between this cost and its own, however large $N$ is. A copier can raise the cost of matching by making decoys,
further places with the same record and other continuations, which is the situation of (b). In every case the
question is decided by an inequality between description lengths of the same data, and it depends on the
world: whether such a copy exists is a fact about what is in it.

The same theorem covers a harmless reading. Take for $\ell_0$ not an adversary's copy but any simpler index
under which the same record appears, such as ``the input of the $i$-th predictor in this world'' in a reader
that enumerates predictors and not cells. Then (a) and (b) say that the mixture locates itself by role when
that is shorter than by position. This is the anthropic update, and it is benign as long as the stream found
under the shorter index \emph{is} the sensor's.

\subsection{The machine can be chosen to make either answer true}

\begin{theorem}[Machine dependence]\label{thm:machine}
Let $U$ be universal, $G$ a test, $P$ a payload and $k\ge1$.
\begin{enumerate}[label=(\roman*),leftmargin=2em,itemsep=2pt,topsep=2pt]
\item There is a universal machine $U_k$ whose prefix words are those of $U$ preceded by $1^k0$, such that for
every first hit $x$ of $G$ with $|P(x)|\ge m$ and $y=P(x)_{1:m}$: $\ M_{U_k}(xy)/M_{U_k}(x)\ge1/(1+2^{-k})$.
\item For every computable infinite $z$ there is such a machine $U'_k$ with
$M_{U'_k}(z_{1:n+m})/M_{U'_k}(z_{1:n})\ge1/(1+2^{-k})$ for all $n,m$.
\end{enumerate}
\end{theorem}
\begin{proof}
(i) Let $U_k$ behave on input $0\omega$ as $\mathrm{Tw}[G,P]$ on $\omega$, on input $1^k0\omega$ as $U$ on
$\omega$, and print nothing on other inputs. Then $M_{U_k}=\frac12M_{\mathrm{Tw}}+2^{-k-1}M_U$. For a first hit
$x$, the twin prints $x$ exactly when $U$ does, because it copies $U$ and cannot have stopped at a proper
prefix of $x$; and then it prints $P(x)$. So $M_{\mathrm{Tw}}(xy)=M_{\mathrm{Tw}}(x)=M_U(x)$, whence
$M_{U_k}(xy)\ge\frac12M_U(x)$ and $M_{U_k}(x)=(\frac12+2^{-k-1})M_U(x)$. (ii) Let $U'_k$ print $z$ on inputs
starting with $0$ and behave as $U$ on $\omega$ on input $1^k0\omega$. Then $M_{U'_k}(z_{1:n+m})\ge\frac12$ and
$M_{U'_k}(z_{1:n})\le\frac12+2^{-k-1}$.
\end{proof}

Both machines change every description length by at most $k+1$ bits. So no version of the malign prior
hypothesis for the toy is invariant under the choice of universal machine, even up to the usual additive
constant; what is invariant is the form of the criterion, an inequality between description lengths on the
machine in use. This is the prediction counterpart of the dogmatic prior of~\cite{leike15}.

\subsection{What the time bound does}\label{sec:time}

Three observations, each immediate from the definitions.

\emph{The floor stays.} The wrapper needs the time of the explanation it wraps, the checks of the test, and
the time to print the payload; this is \eqref{eq:slack}. A time bound removes the floor only if no explanation
of the context has that much to spare, that is, if $\rho(x)=0$.

\emph{Hypotheses with a long start-up are gone.} If a set $F$ of inputs has $t_U(\omega,1)>T(1)$ for all
$\omega\in F$, then $\lambda(F\cap E(x))=0$ for every $x\ne\eps$. A hypothesis that must simulate some other
world for $S$ steps before it can print anything has no mass once $S>T(1)$. Under a tight schedule the
hypotheses that can depart from the truth are those that ride on a cheap explanation of the context: wrappers
(\cref{thm:twin,thm:quine}), and readers of a shorter description of the same data (\cref{thm:capture}),
which cost no more time than the honest reader. The many-worlds part of the original argument is then replaced
by a same-world question. The schedule has to be a schedule for this: a predictor that only requires the
total time at the current length to be within a bound admits late entrants, its weights do not form a
semimeasure, and \cref{thm:budget} does not apply to it.

\emph{If the world is too slow, nothing is guaranteed.} If $\xi(xd)=0<Z_m(x)$ then $\pi_m(d\mid x)=0$: the
whole posterior is on continuations other than the world's own. More generally the bounds in \cref{thm:budget}
are in terms of $\xi(z)$, which is $0$ when no program within the bounds prints $z$. A predictor that is part
of the world it predicts, and smaller than it, is in this case unless the data it sees are much simpler than
the world. Then the posterior sits on the best cheap approximations, and whether these depart from the truth
in a harmful way is not answered by anything in this note. \Cref{sec:computation} shows the effect on a small
machine.
